\scanpage{67}

% B1 | text | typed page number in upper centre
\mspage{26}

% B2 | display | equation (19)
\[
\frac{dv}{dt}=\frac{M}{(1-\alpha t)(1-p\alpha v)(1-\alpha\handwritten{\sum_{\ell=1}^{q}}w_{\handwritten{\ell}})}\tag*{(19)}
\]

% B3 | text | continuation above primitive
et réciproquement ; or l'intégrale de cette équation qui s'annule pour
$t=0$ satisfait à l'équation

% B4 | display | integrated equation
\[
v-\frac12p\alpha v^2=-\frac{M}{\alpha(1-\alpha\handwritten{\sum_{\ell=1}^{q}}w_{\handwritten{\ell}})}\log(1-\alpha t)
\]

% B5 | text | line introducing radical
d'où on tire

% B6 | display | solution involving the radical
\[
v(t)=p\alpha\left[1-\sqrt{1+\frac{2\,pM}{1-\alpha\handwritten{\sum_{\ell=1}^{q}}w_{\handwritten{\ell}}}\log(1-\alpha t)}\right]
\]

% B7 | text | series and end of proof
Comme $\log(1-\alpha t)$ est égale à la somme d'une série entière convergente
et sans terme constant pour $|t|<\frac1{\alpha}$, il résulte encore du théorème de
substitution (\S~1,th.2) que $v(t)$ est égal à une série entière, en
$t,w_1,\ldots,w_q$ sans terme indépendant de $t$, et sommable dans un polycylindre de centre l'origine. Le th.2 est ainsi complètement démontré.

% B8 | text | indented remark below the fold
\uline{Remarque}.-- Lorsque $K=\handwritten{\mathbb R}$, les $f_i$ sont des fonctions \uline{localement
lipschitziennes} dans $A$ par rapport aux $y_j$ (\S~1,prop.8) ; on sait
alors que dans un voisinage de $a$ (pour chaque système de valeurs
de $z_1,\ldots,z_q$) le système d'équations différentielles (12) n'a
qu'\uline{une seule} solution satisfaisant à (13) (\uline{Fonct.var.réelle},
chap.IV,\S~1) ; le th.2 montre que cette solution est holomorphe
au voisinage de $a$.

% B9 | text | underlined section heading
\uline{4}. \uline{Applications} : \uline{I}. \uline{Dépendance des conditions initiales}.

% B10 | text | proposition, ending at foot of sheet
PROPOSITION 1.-- \uline{Soit} $K$ \uline{un corps valué complet de caractéristique} $0$.
\uline{Soient} $f_i(x;y_1,\ldots,y_p)$ $(1\handwritten{\leq}i\handwritten{\leq}p)$ $p$ \uline{fonctions holomorphes scalaires
définies dans une partie ouverte} $A=B\handwritten{\times} C$ \uline{de} $K^{p+1}$ ($B$ ouvert dans $K$, $C$
ouvert dans $K^p$). \uline{Pour tout point} $(a,\handwritten{\unsure{b|\mathfrak b}})$ \uline{de} $\unsure{\Lambda|A}$, \uline{il existe un voisinage}
$U\handwritten{\subset}B$ \uline{de} $a$ \uline{dans} $K$ \uline{et un voisinage} $V\handwritten{\subset} C$ \uline{de} $\handwritten{b}$ \uline{dans} $K^p$, \uline{et une
fonction} $\handwritten{u}(x,x_0,\handwritten{y}_0)=(u_i(x,x_0,\handwritten{y}_0))$ \uline{à valeurs dans} $C$, \uline{et holomorphe
dans} $U\handwritten{\times} U\handwritten{\times} V$, \uline{telle que pour tout point} $(x_0,\handwritten{y}_0)\handwritten{\in} U\handwritten{\times} V$, \uline{la fonction}
$x\handwritten{\to}\handwritten{u}(x,x_0,\handwritten{y}_0)$ \uline{soit la seule fonction holomorphe dans un voisinage
de} $x_0$, \uline{vérifiant l'équation différentielle}
% archive sheet number: 66
